% Part: lambda-calculus
% Chapter: introduction
% Section: lambda-definability

\documentclass[../../../include/open-logic-section]{subfiles}

\begin{document}

\olfileid{lam}{int}{rep}
\olsection{λ-可定义的算术函数}

λ-演算如何能充当计算模型？起初，甚至不清楚该如何理解这一说法。要讨论自然数上的可计算性，我们需要为这些数找到一种合适的表示。下面是一种出奇地有效的表示。

\begin{defn}
对每个自然数~$n$，定义\emph{丘奇数} $\num{n}$ 为 λ-项 $\lambd[x][\lambd[y][(x(x(x(\dots
x(y)))))]]$，其中共有 $n$ 个~$x$。
\end{defn}

项 $\num{n}$ 是“迭代器”：以 $f$ 为输入时，$\num{n}$ 返回将 $y$ 映射为 $f^n(y)$ 的函数。注意每个丘奇数都是范式。现在我们可以界定一个 λ-项“计算”自然数上的函数意味着什么。

\begin{defn}
令 $f(x_0, \dots, x_{k-1})$ 为从 $\Nat$ 到 $\Nat$ 的 $n$ 元偏函数。我们称 λ-项~$F$ \emph{λ-定义}~$f$，当且仅当对每个自然数序列 $n_0$, \dots,~$n_{k-1}$，若 $f(n_0, \dots, n_{k-1})$ 有定义则 \[
F\, \num{n_0}\, \num{n_1} \dots \num{n_{k-1}} \red \num{f(n_0, n_1, \dots,
  n_{k-1})}
\]，否则 $F, \num{n_0}\, \num{n_1}
\dots \num{n_{k-1}}$ 没有范式。
\end{defn}

\begin{thm}
\ollabel{thm:lambda-def}
函数 $f$ 是部分可计算函数，当且仅当它被某个 λ-项所 λ-定义。
\end{thm}

\begin{explain}
这一定理颇令人惊讶。作为一种计算模型，λ-演算是相当简单的演算；仅有的操作就是 λ-抽象与应用！然而，从如此贫乏的资源出发，却能实现任何计算过程。
\end{explain}

\end{document}